Neural networks trained with empirical risk minimization learn spurious correlations, failing systematically on minority groups where spurious features conflict with labels. We demonstrate that simplicity bias creates a detectable signature in training dynamics: minority-group samples exhibit 8$\times$ higher loss at 5\% of training epochs compared to majority samples (Mann-Whitney $p < 10^{-14}$), enabling automatic detection without group labels. Using linear probes on Waterbirds, we verify the mechanism: representations encode spurious features (background) with 91.2\% accuracy versus 80.5\% for core features (bird shape) at epoch 5---a 10.7 percentage point gap persisting throughout training. This loss-based detection achieves 6.6$\times$ improvement over random baseline (33\% vs.\ 5\% precision at matched prediction rate). We provide the first explicit analysis of precision limits under minority imbalance, showing that base rate constraints---not signal weakness---bound individual detection accuracy. Our findings establish early-epoch loss magnitude as a practical signal for minority detection and quantify how simplicity bias manifests as an accuracy magnitude gap (not timing gap) with pretrained features.
